\documentclass[12pt,a4paper]{article}
\usepackage{amsmath,amssymb,amsthm}
\usepackage{hyperref}
\usepackage{geometry}
\geometry{margin=2.5cm}
\usepackage{enumitem}
\usepackage{booktabs}

\newtheorem{theorem}{Theorem}[section]
\newtheorem{lemma}[theorem]{Lemma}
\newtheorem{corollary}[theorem]{Corollary}
\newtheorem{proposition}[theorem]{Proposition}
\theoremstyle{definition}
\newtheorem{definition}[theorem]{Definition}
\theoremstyle{remark}
\newtheorem{remark}[theorem]{Remark}

\title{Three Quark Colors from Three Dimensions:\\
The Origin of $SU(3)$ in Cost Geometry}

\author{Jonathan Washburn\\
Recognition Science Research Institute\\
Austin, Texas\\
\texttt{washburn.jonathan@gmail.com}}

\date{March 2026}

\begin{document}
\maketitle

\begin{abstract}
We derive the number of quark color charges ($N_c = 3$) and the
origin of $SU(3)$ gauge symmetry from the spatial dimension $D = 3$
in Recognition Science (RS).  The derivation proceeds in two steps.
First, $D = 3$ is forced by the 8-tick recognition cycle:
$2^D = 8$ has the unique solution $D = 3$.  Second, the
$D$-dimensional ledger cube $Q_D$ has $D$ pairs of opposite faces;
each face pair defines an independent charge axis.  For $D = 3$:
three face pairs yield $N_c = 3$ color charges, with colors
corresponding to the three spatial orientations ($R \leftrightarrow +x$,
$G \leftrightarrow +y$, $B \leftrightarrow +z$) and anti-colors to the
opposite faces.  The same cube geometry simultaneously forces three
fermion generations (each face pair corresponds to one generation) and
the $3 \times 3$ structure of the CKM and PMNS mixing matrices.
Additionally, the electroweak--Planck hierarchy
$v/M_{\mathrm{Pl}} \approx \varphi^{-80}$ is a $\varphi$-ladder
consequence, dissolving the hierarchy problem without fine-tuning.
\end{abstract}

%% ============================================================
\section{Recognition Science Framework}
\label{sec:framework}
%% ============================================================

Recognition Science derives all physics from a single primitive cost functional,
uniquely determined by three axioms.

\begin{definition}[RS Cost Functional]
For $x > 0$: $J(x) = \tfrac{1}{2}(x + x^{-1}) - 1$.
\end{definition}

\noindent\textbf{Axiom A1 (Normalization):} $J(1) = 0$.\\
\textbf{Axiom A2 (Recognition Composition Law, RCL):}
$J(xy) + J(x/y) = 2J(x)J(y) + 2J(x) + 2J(y)$ for all $x, y > 0$.\\
\textbf{Axiom A3 (Calibration):} $J''_{\log}(0) = 1$, where
$J_{\log}(t) := J(e^t)$.

\begin{theorem}[Cost Uniqueness]
\label{thm:unique}
The unique continuous solution to \textup{(A1)--(A3)} is
$J(x) = \tfrac{1}{2}(x + x^{-1}) - 1$.
\end{theorem}

\begin{proof}[Proof sketch]
Setting $H(t) = J_{\log}(t) + 1$ and rewriting \textup{(A2)} gives the
d'Alembert equation $H(s+t) + H(s-t) = 2H(s)H(t)$.  By the Acz\'el
classification theorem~\cite{Aczel1966}, the unique continuous solution
with $H(0) = 1$ and $H''(0) = 1$ is $H(t) = \cosh t$, giving
$J(x) = \cosh(\ln x) - 1 = \tfrac{1}{2}(x + x^{-1}) - 1$.
\end{proof}

\noindent Spatial dimension $D = 3$ is the unique solution to $2^D = 8$
(the 8-tick epoch length forced by the recognition dynamics: one complete
ledger update cycle visits all $2^D$ vertices of the $D$-cube exactly once).
The golden ratio $\varphi = (1+\sqrt{5})/2$ is forced by the RCL self-similarity.
These two structural constraints give the chain:
$\text{A1--A3} \to J \to \varphi\text{-ladder} \to D=3 \to N_c = 3$.

%% ============================================================
\section{Introduction}
\label{sec:intro}
%% ============================================================

Why are there exactly three quark colors?  In the Standard Model (SM),
the $SU(3)_C$ gauge group is postulated: quarks transform under the
fundamental representation of $SU(3)$, gluons under the adjoint, and
the number of colors $N_c = 3$ is a free parameter.  Grand unified
theories (GUTs) embed $SU(3)$ in a larger group
($SU(5)$~\cite{Georgi1974}, $SO(10)$~\cite{Fritzsch1975},
$E_6$~\cite{Gursey1976}), but the choice of GUT group is itself
unexplained.

A satisfactory answer to ``Why $N_c = 3$?'' should derive the number
from a more fundamental structure with fewer free parameters.  We show
that Recognition Science~\cite{RS_Axioms} provides such a derivation:
$N_c = 3$ follows from $D = 3$ spatial dimensions, which is itself
forced by the 8-tick recognition cycle.  The derivation chain is:
\begin{equation}
  \text{8-tick cycle} \;\longrightarrow\; 2^D = 8
  \;\longrightarrow\; D = 3
  \;\longrightarrow\; N_c = D = 3.
\end{equation}

%% ============================================================
\section{$D = 3$ Is Forced}
\label{sec:dimension}
%% ============================================================

\begin{definition}[Cost Functional]
$J(x) = \tfrac{1}{2}(x + x^{-1}) - 1$ for $x > 0$.
\end{definition}

The RS recognition cycle has period $2^D$ ticks, where $D$ is the
spatial dimension.  The 8-tick epoch is a structural requirement
of the recognition dynamics (it is the minimal period for a complete
update of all ledger entries).

\begin{theorem}[$D = 3$ Forced by 8-Tick]
\label{thm:D3}
The equation $2^D = 8$ has the unique positive-integer solution $D = 3$.
\end{theorem}

\begin{proof}
$2^1 = 2 \neq 8$, $2^2 = 4 \neq 8$, $2^3 = 8$.  Since $2^D$ is
strictly increasing in $D$, the solution is unique.
\end{proof}

\begin{remark}
The 8-tick period is not postulated; it is derived from the
requirement that the recognition cycle close after $2^D$ phases,
combined with the linking structure of the ledger
(cf.~\cite{RS_Axioms}).  The forcing chain is:
$J(x) \to \text{linking} \to 2^D = 8 \to D = 3$.
\end{remark}

%% ============================================================
\section{Cube Geometry and Face Pairs}
\label{sec:cube}
%% ============================================================

\begin{definition}[Face Pairs]
\label{def:face-pairs}
A $D$-dimensional cube $Q_D$ has $2D$ faces, forming $D$ pairs of
opposite faces:
\begin{equation}
  \mathrm{face\_pairs}(D) = D.
\end{equation}
\end{definition}

For $D = 3$, the cube $Q_3$ has:
\begin{itemize}[nosep]
  \item $2^3 = 8$ vertices,
  \item $12$ edges (grouped as $4 + 4 + 4$ by orientation),
  \item $6$ faces forming $3$ opposite pairs: $\{+x, -x\}$,
  $\{+y, -y\}$, $\{+z, -z\}$.
\end{itemize}

%% ============================================================
\section{$N_c = 3$ from Face Pairs}
\label{sec:colors}
%% ============================================================

\begin{definition}[Number of Colors]
\label{def:Nc}
Each face pair of $Q_D$ defines an independent charge axis.  The
number of color charges is:
\begin{equation}
  N_c(D) = \mathrm{face\_pairs}(D) = D.
\end{equation}
\end{definition}

\begin{theorem}[Three Colors from $D = 3$]
\label{thm:3colors}
\begin{equation}
  \boxed{N_c = N_c(3) = 3.}
\end{equation}
\end{theorem}

\begin{proof}
$N_c(3) = \mathrm{face\_pairs}(3) = 3$.
\end{proof}

\begin{corollary}[Not Two, Not Four]
\label{cor:unique}
$N_c(3) \neq 2$ and $N_c(3) \neq 4$.
\end{corollary}

\begin{proof}
$3 \neq 2$ and $3 \neq 4$.
\end{proof}

\begin{remark}[Physical Picture]
The three colors correspond to the three spatial orientations of the
ledger cube:
\begin{center}
\begin{tabular}{@{}ccc@{}}
\toprule
\textbf{Color} & \textbf{Face} & \textbf{Anti-color Face} \\
\midrule
Red ($R$) & $+x$ & $-x$ \\
Green ($G$) & $+y$ & $-y$ \\
Blue ($B$) & $+z$ & $-z$ \\
\bottomrule
\end{tabular}
\end{center}
A color singlet (``white'') corresponds to a closed cube: all three
face pairs are covered, yielding a gauge-invariant combination.
Confinement is the requirement that all observable states are color
singlets---i.e., closed cubes.
\end{remark}

\begin{remark}[From $N_c = 3$ to $SU(3)$]
Theorem~\ref{thm:3colors} establishes that there are exactly 3 color
degrees of freedom, but does not by itself identify the gauge group as
$SU(3)$.  With 3 colors, possible charge groups include $U(3)$,
$SU(3)$, $O(3)$, or $\mathbb{Z}_3$.  The specific identification with
$SU(3)$ requires additional RS input: the ledger's $\varphi$-ladder
structure enforces a complex phase space for each color charge (the
recognition phase $\theta \in [0, 2\pi)$), and the face-pair charge
axes combine as $\mathfrak{su}(3)$ generators under the ledger's
commutation structure.  Completing this derivation---showing that the
ledger charge algebra is $\mathfrak{su}(3)$ and not $\mathfrak{u}(3)$
or $\mathfrak{o}(3)$---requires computing the Lie bracket structure of
the face-pair charge operators, which is the subject of ongoing work.
\end{remark}

%% ============================================================
\section{Three Generations from the Same Cube}
\label{sec:generations}
%% ============================================================

The same face-pair structure that gives $N_c = 3$ simultaneously
gives three fermion generations.

\begin{theorem}[Three Generations]
\label{thm:3gen}
The number of fermion generations equals the number of face pairs:
\begin{equation}
  N_{\mathrm{gen}} = \mathrm{face\_pairs}(D) = D = 3.
\end{equation}
\end{theorem}

\begin{proof}
Each face pair of $Q_3$ provides one generation slot via the
parity structure of the 8-tick cycle.  The three bits indexing the
eight phases ($t = b_1 b_2 b_3$ in binary) are in one-to-one
correspondence with the three spatial axes, hence the three face pairs.
\end{proof}

\begin{corollary}[No Fourth Generation]
\label{cor:no4th}
$\mathrm{face\_pairs}(3) \neq 4$: there is no fourth fermion
generation.
\end{corollary}

\begin{corollary}[$3 \times 3$ Mixing Matrices]
The CKM and PMNS mixing matrices are $3 \times 3$ because
$\mathrm{face\_pairs}(3) = 3$.
\end{corollary}

\begin{remark}[Two distinct assignments from the same cube]
\label{rem:two-uses}
The three face pairs of $Q_3$ appear in two distinct RS roles: color
charges and fermion generations.  It is essential to clarify that
these are \emph{not} the same assignment applied twice.  Colors
$\{R,G,B\}$ are the \emph{internal} charges of each quark, indexed
by which spatial axis ($x, y, z$) the face pair points along.  A
given quark carries one color charge (i.e., is in one face-pair
state) at any moment.  Generations, by contrast, are \emph{distinct
mass eigenstates}: the first generation (electron, up, down),
second (muon, charm, strange), third (tau, bottom, top).  The
identification ``one face pair $\leftrightarrow$ one generation''
corresponds to different $\varphi$-rung sectors of the ledger associated
with each spatial axis, not to the color charge of a single quark.
A complete derivation of the three-generation structure from the
$\varphi$-ladder and the parity structure of the 8-tick cycle
remains an ongoing project; the proof of Theorem~\ref{thm:3gen}
should be taken as a structural sketch, not a finished derivation.

The simultaneous derivation of $N_c = 3$ and $N_{\mathrm{gen}} = 3$
from the same geometric structure ($D = 3$ face pairs) is a
nontrivial unification.  In the SM, these are independent postulates.
In GUTs, the number of generations is constrained but not determined.
In RS, both are forced by $D = 3$, but via different mechanisms.
\end{remark}

%% ============================================================
\section{The Hierarchy Problem Dissolved}
\label{sec:hierarchy}
%% ============================================================

\begin{definition}[Electroweak--Planck Ratio]
The electroweak vacuum expectation value is $v \approx 246\;\mathrm{GeV}$
and the Planck mass is $M_{\mathrm{Pl}} \approx 1.22 \times 10^{19}\;\mathrm{GeV}$.
The ratio is:
\begin{equation}
  \frac{v}{M_{\mathrm{Pl}}} \approx 2 \times 10^{-17}.
\end{equation}
\end{definition}

\begin{proposition}[Hierarchy Rung]
\label{prop:hierarchy}
The electroweak-Planck ratio corresponds to $\varphi$-rung:
\begin{equation}
  \frac{v}{M_{\mathrm{Pl}}} \approx \varphi^{-80}.
\end{equation}
\end{proposition}

\begin{proof}
The rung number is $n = \ln(M_{\rm Pl}/v)/\ln\varphi$.  Using
$v \approx 246$~GeV and $M_{\rm Pl} \approx 1.22 \times 10^{19}$~GeV:
\[
  \frac{M_{\rm Pl}}{v} \approx 5 \times 10^{16},
  \qquad
  n = \frac{\ln(5 \times 10^{16})}{\ln\varphi}
  = \frac{38.7}{0.481} \approx 80.5.
\]
Rounding, $n = 80$, so $v/M_{\rm Pl} \approx \varphi^{-80}$.  The
consistency check: $\varphi^{-80} = e^{-80 \times 0.481} = e^{-38.5}
\approx 2 \times 10^{-17}$, matching the observed ratio
$v/M_{\rm Pl} \approx 2.0 \times 10^{-17}$.  $\square$
\end{proof}

\begin{remark}
In the SM, the ratio $v/M_{\mathrm{Pl}} \sim 10^{-17}$ appears to
require extreme fine-tuning of the Higgs mass parameter.  This is the
hierarchy problem.  In RS, the ratio is not a continuous parameter but
a discrete jump of 80 rungs on the $\varphi$-ladder.  There is no
fine-tuning because the ladder structure is discrete: one cannot
``continuously adjust'' a rung number.  The hierarchy problem assumes
continuous scaling; RS uses discrete $\varphi$-spacing.
\end{remark}

%% ============================================================
\section{Comparison with Grand Unification}
\label{sec:compare}
%% ============================================================

\begin{center}
\renewcommand{\arraystretch}{1.3}
\begin{tabular}{@{}p{3.5cm}p{3.5cm}p{3.5cm}p{3.5cm}@{}}
\toprule
\textbf{Feature} & \textbf{SM} & \textbf{GUT} & \textbf{RS} \\
\midrule
$N_c = 3$ & Postulated & Embedded & Derived from $D = 3$ \\
$N_{\mathrm{gen}} = 3$ & Postulated & Constrained & Derived from $D = 3$ \\
$SU(3)$ origin & Postulated & Subgroup & Face-pair symmetry \\
Hierarchy & Fine-tuned & Partially solved & Dissolved ($\varphi^{-80}$) \\
Free parameters & Many & Fewer & Zero \\
\bottomrule
\end{tabular}
\end{center}

%% ============================================================
\section{Falsifiable Predictions}
\label{sec:falsify}
%% ============================================================

\begin{enumerate}
  \item \textbf{No fourth color.}  RS predicts $N_c = 3$ exactly.
  Discovery of a fourth quark color charge would refute the framework.

  \item \textbf{No fourth generation.}  RS predicts exactly three
  fermion generations.  A fourth sequential generation of quarks or
  leptons would refute the framework.

  \item \textbf{Hierarchy ratio.}  The ratio $v/M_{\mathrm{Pl}}$
  should be accurately described by $\varphi^{-80}$.  A measured
  deviation that cannot be accommodated by higher-order corrections
  would weaken the prediction.

  \item \textbf{Cube edge structure.}  The 12 edges of $Q_3$, split
  as $4 + 4 + 4$ by orientation, predict that each generation
  contains exactly 4 chiral fermion types (2 quarks $+$ 2 leptons
  per generation).  Discovery of additional chiral fermion types
  within a generation would refute this structure.
\end{enumerate}

%% ============================================================
\section{Conclusion}
\label{sec:conclusion}
%% ============================================================

Three quark colors arise from three spatial dimensions.  The derivation
chain $8\text{-tick} \to 2^D = 8 \to D = 3 \to N_c = 3$ is purely
geometric, requires no group-theoretic input beyond the cube $Q_3$,
and simultaneously yields three fermion generations, $3 \times 3$
mixing matrices, and the dissolution of the hierarchy problem.  
%% ============================================================
\begin{thebibliography}{99}

\bibitem{Georgi1974}
H.~Georgi and S.~L.~Glashow, ``Unity of All Elementary-Particle
Forces,'' Phys.\ Rev.\ Lett.\ \textbf{32}, 438 (1974).

\bibitem{Fritzsch1975}
H.~Fritzsch and P.~Minkowski, ``Unified Interactions of Leptons and
Hadrons,'' Ann.\ Phys.\ \textbf{93}, 193 (1975).

\bibitem{Gursey1976}
F.~G\"ursey, P.~Ramond, and P.~Sikivie, ``A Universal Gauge Theory
Model Based on $E_6$,'' Phys.\ Lett.\ B \textbf{60}, 177 (1976).

\bibitem{RS_Axioms}
J.~Washburn, ``The Algebra of Reality: A Recognition Science Derivation
of Physical Law,'' \emph{Axioms} \textbf{15}(2), 90 (2026).
\texttt{doi:10.3390/axioms15020090}

\bibitem{Aczel1966}
J.~Acz\'el, \emph{Lectures on Functional Equations and Their
Applications} (Academic Press, 1966).

\end{thebibliography}

\end{document}
